\documentclass[12pt]{article}

\usepackage[T1]{fontenc}
\usepackage[utf8]{inputenc}
\usepackage{geometry}
\usepackage{setspace}
\usepackage{parskip}
\usepackage{hyperref}

\geometry{margin=1in}
\onehalfspacing

\title{The Last Remaining Asymmetry:\\Fairness, Competition, and the Future of Fertility}
\author{Kangbin Yim\thanks{Email: yim@sch.ac.kr (Prof., SCH University, Korea)}}
\date{May 2026}


\begin{document}

\maketitle

For decades, fertility decline has been explained primarily through economics.
Housing costs are too high. Education is too expensive. Employment is unstable.
The opportunity cost of raising children continues to increase. 

Governments around the world have responded with subsidies, tax incentives, childcare programs, and various forms of family support. Yet despite these efforts, fertility rates continue to decline across many highly developed societies. Perhaps this suggests that the problem is not purely economic. Perhaps something deeper is occurring beneath the surface. 

Most discussions about fertility focus on cost. Few focus on the structure of civilization itself. Yet throughout history, civilization has followed a remarkably consistent direction.

It has gradually reduced asymmetry. The distance between rulers and citizens has narrowed. The distinction between aristocrats and commoners has weakened.
Educational opportunities have become more accessible. Political participation has become more inclusive. Economic participation has become less dependent upon inherited privilege.

In many ways, the history of civilization can be understood as a long process of expanding fairness. But fairness produces a consequence that is often overlooked. Fairness does not eliminate competition. It universalizes competition.

In a highly stratified society, many individuals are excluded from competition before it even begins. In a more symmetrical society, more people enter the same arena. The result is not less competition, but broader competition. The modern world is therefore not merely a fairer world. It is also a more competitive world.

This transformation is especially visible in the relationship between men and women. For most of human history, men and women occupied substantially different social roles. The reasons were not purely cultural. They were also functional.

Survival depended heavily upon physical strength, environmental risk, warfare, hunting, and labor-intensive production. Men and women participated in society through different pathways because society itself was organized around different requirements. 

Whether these arrangements were fair or unfair is a separate question. The important point is that the competition itself was often separated. Men and women were frequently competing within different domains.

Modern civilization has steadily dismantled these distinctions. Industrialization reduced the importance of physical strength. Mass education expanded access to knowledge. Digital technology reduced barriers to information. Professional opportunities expanded. Economic independence expanded. Political rights expanded.

As a result, men and women increasingly participate within the same systems of competition. They attend the same schools. Compete for the same university admissions. Compete for the same professional positions. Compete for the same promotions. Compete for the same leadership roles. Compete for the same definitions of success. The trajectory is clear. The modern world is not merely creating greater equality. It is creating a shared competitive arena. 

Today, artificial intelligence and robotics are accelerating this process. Machines have already replaced much of the physical labor that once favored male strength. Artificial intelligence is beginning to automate forms of cognitive labor that once created other forms of differentiation.

Robots increasingly perform physically demanding work. Algorithms increasingly perform analytical work. As these technologies advance, biological sex becomes progressively less relevant to economic productivity and professional achievement. The long-term effect may be unprecedented. 

For the first time in human history, men and women may find themselves competing within almost entirely symmetrical social structures. And it is precisely at this moment that a paradox emerges. As more asymmetries disappear, one asymmetry becomes increasingly visible. It's the pregnancy and childbirth. 

Pregnancy occurs within the female body. Childbirth occurs within the female body. No political reform can eliminate this reality. No economic policy can completely redistribute it. No legal framework can fully equalize it. It remains a biological fact. Paradoxically, the more symmetrical society becomes, the more visible this remaining asymmetry appears. 

In earlier societies, asymmetries existed everywhere. Men often carried greater physical risk through warfare, hunting, dangerous labor, and external threats. Women carried greater responsibility for pregnancy, childbirth, and childrearing. Society itself was structured around multiple asymmetries. Because asymmetry was widespread, no single asymmetry dominated public consciousness. 

Today, however, many of those asymmetries have weakened. Educational asymmetries have weakened. Political asymmetries have weakened. Professional asymmetries have weakened. Economic asymmetries have weakened. As these differences fade, childbirth increasingly stands alone. It becomes one of the few remaining asymmetries that cannot be fully removed from the system. Yet the significance of this asymmetry extends beyond biology.

In a society organized around fair competition, childbirth may increasingly be perceived as an asymmetric burden within a symmetrical arena. This does not imply that women consciously reject childbirth because they regard it as unfair. Human decisions are far more complex than that. Economic conditions matter. Cultural values matter. Personal aspirations matter. Public policy matters. The argument is more subtle.

It is that the meaning of childbirth may change when viewed against the background of expanding fairness and expanding competition. An asymmetry that once existed among many asymmetries now exists among few. An asymmetry that once existed within separate social pathways now exists within shared competitive structures. As a result, it becomes increasingly visible. And what becomes increasingly visible may also become increasingly consequential.

This perspective may help explain why fertility decline often appears most severe in highly developed societies. Advanced societies are not merely wealthier. They are more symmetrical. They are more competitive. They provide greater educational opportunity. Greater professional mobility. Greater individual autonomy. And increasingly, they place men and women within the same systems of evaluation and competition.

The very forces that expand freedom and fairness may simultaneously increase awareness of biological asymmetry. From this perspective, fertility decline may not be merely a demographic phenomenon. It may represent a tension between two historical trajectories. 

The first is civilization’s movement toward greater fairness and broader competition. The second is the biological structure of human reproduction. For most of history, these trajectories coexisted without major conflict because society itself contained numerous asymmetries. As civilization removes more of those asymmetries, reproduction becomes increasingly exceptional.

The question facing advanced societies may therefore be different from the one policymakers usually ask. The question may not simply be: “How can we reduce the economic cost of having children?” A deeper question may be: “What happens when a society built upon fairness and universal competition encounters a biological asymmetry that cannot be fully removed?”

Artificial intelligence and robotics may make this question even more important. As technology continues to erase distinctions in labor, capability, and social participation, the visibility of biological asymmetry may continue to increase. The future challenge of fertility may therefore be less about economics alone and more about how civilization responds to the last remaining asymmetry.

\end{document}
